% ===================================================================
%  اللوحة رقم ٢ — جسم الإنسان. — الفسحة. — الألعاب.
%  Traduction 2026 d'apres texto/io, controlee sur texto/fr, texto/en
%  et texto/es. Consigne : voir texto/ar/10-lawha-01.tex.
% ===================================================================

%%K t02-apar-1 apar
\VUsaut{4.0mm}
\VUtitre{15.0pt}{60}{اللوحة رقم ٢}
\VUsaut{2.0mm}
\VUtitre{11.4pt}{0}{جسم الإنسان. --- الفسحة.}
\VUtitre{11.4pt}{0}{الألعاب.}

%%K t02-tit-0 sub
\VUtitre{13.2pt}{0}{جسم الإنسان.}
\VUsaut{2.4mm}
\VUcentre{10.2pt}{0}{\textit{(درس بسيط في العلوم الطبيعية.)}}
\VUsaut{3.0mm}

%%K t02-01-1 p
1. --- أجزاء جسم الإنسان الرئيسية هي: \VUgras{الرأس} و\VUgras{الجذع}
و\VUgras{الأطراف}.

%%K t02-tit-1 sub
\VUpk{10.2pt}{40}{الرأس.}
\VUsaut{3.0mm}

%%K t02-02-1 p
2. --- للرأس جزءان: \VUgras{الجمجمة} التي تحوي \VUgras{الدماغ} ويغطيها
\VUgras{الشعر}، و\VUgras{الوجه}.

%%K t02-03-1 p
3. --- في رسمنا لا نرى الجمجمة ولا الدماغ؛ لكننا نرى أجزاء الوجه المهمة
رؤية واضحة جدًا.

%%K t02-04-1 p
4. --- فهذه أولًا \VUgras{الجبهة}، وهي تنم عن ذكاء قوي وفكر لا يكف عن
العمل. و\VUgras{الصدغان} منفرجان جدًا. و\VUgras{العين} حية مفتوحة على
سعتها. يحميها \VUgras{حاجب} كثيف و\VUgras{أهداب} طويلة.

%%K t02-05-1 p
5. --- وهنا أيضًا \VUgras{الأنف} و\VUgras{المنخران}. الأنف يخدمنا في الشم
وفي التنفس. وعلى جانبي الأنف \VUgras{الخدان}، وأبعد منهما \VUgras{الأذنان}.
ويتألف أسفل الوجه من \VUgras{الفم} و\VUgras{الذقن}.

%%K t02-06-1 p
6. --- ولا نراهما جيدًا، لأن \VUgras{الشارب} و\VUgras{اللحية} يخفيانهما
عنا. لكننا نعلم أن في الفم \VUgras{الشفتين} و\VUgras{الفك الأعلى}
و\VUgras{الفك الأسفل} بما فيهما من \VUgras{أسنان}، و\VUgras{اللسان}. بالفم
نتكلم ونأكل. وباللسان نتذوق طعامنا.

%%K t02-07-1 p
7. --- والعين والأذن والأنف والفم، مع الأصابع، هي أعضاء الحواس الخمس:
\VUgras{البصر} و\VUgras{السمع} و\VUgras{الشم} و\VUgras{الذوق}
و\VUgras{اللمس}.

%%K t02-08-1 p
8. --- فالرأس إذًا من أطرف أجزاء جسم الإنسان.

%%K t02-tit-2 sub
\VUsaut{4.4mm}
\VUpk{10.2pt}{40}{الجذع.}
\VUsaut{3.0mm}

%%K t02-09-1 p
9. --- \VUgras{العنق} يصل الرأس بالجذع. ويشمل \VUgras{الحلق}
و\VUgras{القفا}.

%%K t02-10-1 p
10. --- والجذع يحوي كذلك أعضاء كثيرة لا غنى عنها. ففي \VUgras{الصدر}
\VUgras{الرئتان} اللتان نتنفس بهما، و\VUgras{القلب} الذي هو مركز الحياة.
فإذا كف القلب عن الخفقان مات الكائن الحي.

%%K t02-11-1 p
11. --- و\VUgras{الأضلاع} تسند الصدر.

%%K t02-12-1 p
12. --- وسائر أجزاء الجذع هي \VUgras{الخاصرة} و\VUgras{الظهر}
و\VUgras{الكتفان} و\VUgras{البطن}. ننام على خاصرتنا أو على ظهرنا، ونحمل
الأثقال على كتفنا.

%%K t02-13-1 p
13. --- وفي البطن \VUgras{المعدة} و\VUgras{الأمعاء}. وهما تخدماننا في هضم
طعامنا.

%%K t02-tit-3 sub
\VUsaut{4.4mm}
\VUpk{10.2pt}{40}{الأطراف.}
\VUsaut{3.0mm}

%%K t02-14-1 p
14. --- لنا أربعة \VUgras{أطراف}: \VUgras{ذراعان} و\VUgras{ساقان}.
بذراعينا نأخذ الأشياء ونمسكها ونرفعها ونجمعها؛ وبساقينا نقف ونمشي ونجري.

%%K t02-15-1 p
15. --- \VUgras{الكتف} يصل الذراع بالجسم. و\VUgras{عضلة} متينة جدًا، هي
ذات الرأسين، تحرّك الذراع التي يصلها \VUgras{المرفق} بـ\VUgras{الساعد}.
و\VUgras{الرسغ} هو مفصل \VUgras{اليد} التي تنتهي بـ\VUgras{الأصابع}
و\VUgras{الأظافر}. وعلى اليد يتبين \VUgras{وريد} (وشريان) يجري فيه
\VUgras{الدم}.

%%K t02-16-1 p
16. --- وأجزاء الساق هي: \VUgras{الفخذ} و\VUgras{الركبة}
و\VUgras{مأبضها}، و\VUgras{بطة الساق} التي يغطي \VUgras{لحمها}
\VUgras{الجلد}، و\VUgras{الكاحل} و\VUgras{القدم}. و\VUgras{عظام} الساق
غليظة جدًا وقوية جدًا. وفي طرف القدم \VUgras{أصابع القدم}. وسائر الأجزاء
هي \VUgras{الأخمص} و\VUgras{العقب}.

%%K t02-17-1 p
17. --- هذا وصف جسم الإنسان وصفًا يكاد يكون تامًا.

%%K t02-17-2 p
\textit{ملاحظة.} --- \textit{بعبارة «يؤلمني... (رأسي، وهكذا)» يستطيع
المعلم أن يعيد هذه الألفاظ كلها من جهة} \VUgras{صحة الجسم}
\textit{أو اعتلالها.}

%%K t02-tit-4 sub
\VUsaut{5.0mm}
\VUpk{10.2pt}{40}{الرياضة البدنية.}
\VUsaut{2.4mm}
\VUcentre{10.2pt}{0}{\textit{(يرويها أحد التلاميذ.)}}
\VUsaut{3.0mm}

%%K t02-18-1 p
18. --- لكي نكون لِينين خِفافًا أصحّاء، علينا أن نمرّن سيقاننا وأذرعنا.
ولذلك نلعب ونمارس \VUgras{الرياضة البدنية}.

%%K t02-19-1 p
19. --- في أيام الأسبوع يجري هذان التمرينان في فناء المدرسة (انظر اللوحة
رقم ١). أما يومي الخميس والأحد فنخرج إلى مرج واسع، فيه متسع للجري
واللهو.

%%K t02-20-1 p
20. --- ومعلمونا وآباؤنا يأتون معنا أحيانًا؛ لكن الذي يدير التمارين هو
\VUgras{معلم الرياضة} \textsuperscript{(12)}.

%%K t02-21-1 p
21. --- في المرج، على يسار المدخل، تقوم \VUgras{أجهزة الرياضة}
\textsuperscript{(1)}: نصعد \VUgras{السلّم} \textsuperscript{(2)}، ونتأرجح منكوسين
على \VUgras{الحلقتين} \textsuperscript{(3)} وعلى \VUgras{الأرجوحة المعلقة}
\textsuperscript{(4)}، ونرتقي بأيدينا إلى أعلى \VUgras{سلّم الحبال}
\textsuperscript{(5)} أو \VUgras{الحبل المعقود} \textsuperscript{(6)}.

%%K t02-22-1 p
22. --- وفي هذه الأثناء يقفز رفاقنا فوق \VUgras{الحصان الخشبي}
\textsuperscript{(7)} أو \VUgras{المتوازيين} \textsuperscript{(11)}. وآخرون
\VUgras{يلاكمون} \textsuperscript{(8)} أو \VUgras{يبارزون} \textsuperscript{(9)} بقناع
و\VUgras{شيش}. وآخرون أخيرًا يرفعون \VUgras{الأثقال} \textsuperscript{(10)}.

%%K t02-tit-5 sub
\VUsaut{4.6mm}
\VUpk{10.2pt}{40}{الألعاب.}
\VUsaut{3.0mm}

%%K t02-23-1 p
23. --- وحول اللاعبين يلعب الأولاد والبنات \VUgras{ألعابًا} شتى. فالبنات
يلهون بـ\VUgras{الطوق} \textsuperscript{(14)}؛ ويلعبن \VUgras{نطّ الحبل}
\textsuperscript{(13)}، أو يدعون إخوتهن وأولاد عمومتهن إلى دور من
\VUgras{الكروكيه} \textsuperscript{(15)}. وما أسعدهن حين يقصين \VUgras{كرة}
الخصم بـ\VUgras{مطرقتهن الخشبية} وهو يحسب أنه سيمر من القوس بسهولة!

%%K t02-24-1 p
24. --- أما الأولاد فيؤثرون الألعاب الأشدّ. يلعبون عن طيب خاطر
\VUgras{الأوتاد} \textsuperscript{(16)}، يسقطونها بـ\VUgras{الكرة}
\textsuperscript{(17)} إسقاطًا محكمًا. ولا يزدرون \VUgras{الخذروف}
\textsuperscript{(18)} و\VUgras{الكأس والكرة} \textsuperscript{(19)}، ولا حتى
\VUgras{لعبة الضفدع} \textsuperscript{(20)}. لكن لعبتيهم المفضلتين هما
\VUgras{الكُرات} \textsuperscript{(21)} و\VUgras{كرة الجدار} \textsuperscript{(22)}،
فجدار \VUgras{الدفيئة} \textsuperscript{(26)} يصلح تمامًا لِرَدّ الكرة الخفيفة.

%%K t02-25-1 p
25. --- ويوحنا الصغير، فوق \VUgras{ركائزه} \textsuperscript{(24)}، ماهر جدًا
ويحفظ توازنه أحسن حفظ. وقربه بضعة رفاق يلعبون \VUgras{الغميضة}
\textsuperscript{(25)} في تلك الدفيئة وخلف \VUgras{السياج}. وآخرون يؤثرون
\VUgras{لعبة الكف} \textsuperscript{(28)} أو \VUgras{الريشة} \textsuperscript{(29)}
التي تطير من \VUgras{مضرب} إلى آخر.

%%K t02-noto-1 noto
\VUnotes{143.2mm}{%
(*) كثيرًا ما تحمل ألعاب الأطفال أسماء قومية بحتة. وقد سميناها بالإيدو
على قدر وسعنا: تارة بالحركة التي تميزها، وتارة بأخذ اسمها القومي في هذا
البلد أو ذاك إذا لم يكن مضلِّلًا أكثر مما ينبغي. مثال ذلك:
\VUgras{لعبة البرميل} (الألمانية والفرنسية) هي \VUgras{لعبة الضفدع}
\textsuperscript{(20)} (الإنجليزية)، وفيها يحاول اللاعبون أن يقذفوا أقراصهم
المستديرة في فم الضفدع المعدني الفاغر فوق الصندوق المثقوب (البرميل).
و\VUgras{قفزة الكتفين} \textsuperscript{(30)} لا تختلف عن \VUgras{قفزة الظهر}
إلا في هذا: أن اللاعبين، وهم يقفزون فوق الرأس، يستندون بأيديهم إلى كتفي
الرفيق، بينما في «\VUgras{قفزة الظهر}» لا يستندون إلا إلى ظهره. --- أو
كذلك: ينحني بعض اللاعبين ويثب الآخرون فوق ظهورهم مستندين بأيديهم إلى
أكتافهم؛ فعلى الأولين أن يحتملوا الصدمة دون أن ينهاروا، وعلى الآخرين أن
يثبوا دون أن يفقدوا توازنهم (انظر اللوحة نفسها).%
}

%%K t02-26-1 p
26. --- وفي وسط المرج شجرتان. وعند أصل إحداهما أقام سبعة أو ثمانية أولاد
لعبة \VUgras{قفزة الكتفين} \textsuperscript{(30)} (*). وهذه اللعبة ليست معروفة
في كل البلدان؛ لكن الجميع يعرفون ما \VUgras{قفزة الظهر}، وهي التي لا
يلعبها الأطفال الآن.

%%K t02-tit-6 sub
\VUsaut{5.0mm}
\VUpk{10.2pt}{40}{الألعاب في الهواء الطلق.}
\VUsaut{3.4mm}

%%K t02-27-1 p
27. --- و\VUgras{الاستغماية} \textsuperscript{(31)} كذلك مسلية جدًا؛ تتناوب
البنات والأولاد على وضع \VUgras{العصابة} على العيون وعلى الإمساك
بالمتثاقلين الذين يَدَعون أنفسهم يُمسَكون.

%%K t02-28-1 p
28. --- وفي وسط المرج، ها هنا لعبة فرنسية جدًا، وإن تكن خشنة بعض الشيء:
إنها \VUgras{الدب} \textsuperscript{(32)}، وهي أشد ضجيجًا بكثير من لعبة
\VUgras{الطائرة الورقية} \textsuperscript{(33)} الهادئة، بل من لعبة
\VUgras{التنس} \textsuperscript{(34)} الإنجليزية.

%%K t02-29-1 p
29. --- وهذه الرياضة الأخيرة متعة التلاميذ الكبار. ومعلمونا أنفسهم
يقبلون عليها عن طيب خاطر. وهي تقتضي مهارة كبيرة وخفة، وأنا أحبها كثيرًا.
أما \VUgras{الأرجوحة} \textsuperscript{(35)} فأتركها للبنات.

%%K t02-30-1 p
30. --- ولا تروقني كذلك لعبة إنجليزية أخرى قد تصير خطرة.

%%K t02-30-1b p
أعني \VUgras{كرة القدم} \textsuperscript{(36)}، وهي فرحة أخي الأكبر. وأنا
أوثر لعبتنا القديمة \VUgras{الأسرى} \textsuperscript{(38)}، وهي نافعة للبدن
مثلها وأقل منها عنفًا بكثير.

%%K t02-tit-7 sub
\VUsaut{4.6mm}
\VUpk{10.2pt}{40}{الدراجة وألعاب الكرة.}
\VUsaut{3.0mm}

%%K t02-31-1 p
31. --- ويسرني أيضًا أن أطوف بالمرج على \VUgras{دراجة} \textsuperscript{(37)}.

%%K t02-32-1 p
32. --- وفي السنة المقبلة سأتعلم \VUgras{ركوب الخيل} \textsuperscript{(39)}. ما
أجمل الحصان حين يخطر أو يهرول في رشاقة، ويجتاز الحواجز عَدْوًا!

%%K t02-33-1 p
33. --- وفي الصيف نتعلم \VUgras{السباحة} \textsuperscript{(40)}. نغوص ونستحم
مستمتعين في ماء النهر الصافي البارد. وأحيانًا نطلق في الختام
\VUgras{بالونًا ورقيًا} \textsuperscript{(41)} يرتفع في الهواء في جلال ما إن
يمتلئ بالهواء الساخن.

%%K t02-34-1 p
34. --- وثمة ألعاب أخرى كثيرة، لكني لن أفرغ من عدّها أبدًا، ويُرى من هذا
الحديث القصير أن أيام الخميس على مرجنا الجميل ليست مملة بحال.

%%K t02-34-2 p
N.~B. --- \textit{يسهل على المعلم أن يتم هذا الفصل بوصف كل لعبة من
الألعاب الأهم وصفًا أدق.}
\VUsaut{6.0mm}
\VUfilet{22mm}
